%%%%%%%%%%%%%%%%%%%%%%%%%%%%%%%%%
% Introduction                  %
%%%%%%%%%%%%%%%%%%%%%%%%%%%%%%%%%

\section{Introduction}

This chapter introduces the host company, EY, and the practice in which the work was carried out, describes the banking and regulatory context that motivates the study, and states the problem. It then defines the purpose, the research objectives, and the research questions, delimits the scope and boundaries of the work, summarizes the contributions, and outlines the structure of the report.

\subsection{Host Company}

The project was carried out as a six-month internship within EY.

\subsubsection{Company Introduction}

Ernst \& Young, known as EY, is one of the four largest professional services firms in the world, alongside Deloitte, PwC, and KPMG, the group known as the Big Four. It operates as a global network of member firms across four service lines: Assurance, Tax, Strategy and Transactions, and Consulting, and positions data and technology as central levers of the value it delivers to clients. EY Tunisia, the local member firm, delivers consulting, assurance, and tax services to leading organizations in the country and has a long standing presence in the financial sector, where banks, insurers, and regulators form a significant share of its client base.

This project was conducted within the Technology Strategy \& Transformation (TST) practice of EY Tunisia, part of the Consulting service line, which accompanies clients through technology transformation with a focus on measurable business outcomes. Its work spans three capability areas: technology strategy, which aligns technology investments with business objectives; next-generation technology operations; and infrastructure and service resilience, including cloud migration across datacenters, networks, and storage. The team counts some thirty-five practitioners in Tunis, led by a partner, and has delivered more than one hundred and ten projects over the past decade, a large share of them offshore for clients across Africa, the MENA region, and Europe. Its sector coverage spans government and the public sector, private enterprises, technology, media and telecommunications, and financial services, the sector of the present project. In banking specifically, the practice combines technology-strategy expertise with regulatory awareness, precisely the intersection at which this project sits. The internship therefore offered direct access to the methodological assets of the EY network, to practitioners experienced with the Tunisian financial sector, and to the regulatory concerns that shape data programs in it.

\subsubsection{Goals and Objectives of the Company}

The Technology Strategy \& Transformation practice supports organizations, among them banks and financial institutions, in their technology and digital transformation, regulatory compliance, and data initiatives. It combines local knowledge of the Tunisian market with the global methodologies of the EY network to help institutions align their technology, and increasingly their data, with their business objectives. A recurring difficulty is that clients ask a question the practice cannot yet answer in a standardized way: where do we stand on data maturity, and how do we compare with our peers and with what the regulator now expects. Data maturity here means the degree to which an organization governs, controls, and exploits its data as a strategic asset. Answering this consistently requires a shared measurement instrument, grounded in the banking context, aligned with the Tunisian regulatory environment, and repeatable across institutions and over time. This is the context of the present project, whose objective was to design a data maturity assessment framework for the Tunisian banking sector and an interactive tool that operationalizes it, moving the practice from qualitative judgment toward an evidence based and auditable diagnosis.

\subsection{Project Context}

The Tunisian banking sector is undergoing a fundamental transformation driven by two convergent forces: the acceleration of digital finance and a tightening regulatory landscape. On the demand side, the spread of mobile banking, instant payments, and data driven services raises customer expectations and makes data reliability a competitive rather than purely technical concern. On the supply side, banks accumulate data across core banking systems, credit and risk platforms, and channel applications rarely designed to interoperate, leaving it fragmented and hard to trust end to end.

On the regulatory front, the Banque Centrale de Tunisie (BCT) issued Circulaire N\textdegree{}2025-08\footnote{Banque Centrale de Tunisie, Circulaire N\textdegree{}2025-08 concerning data governance obligations, published 16~May 2025; its data-governance obligations are set out in Article~6 of the circular. Available at \url{https://www.bct.gov.tn}.}, published on 16 May 2025, mandating structured data governance across all licensed banking institutions. It addresses accountability for data, the formalization of data policies, the classification and retention of information, and the availability of documentary evidence for supervision. In parallel, the internationally recognized BCBS 239\footnote{Basel Committee on Banking Supervision, \emph{Principles for Effective Risk Data Aggregation and Risk Reporting} (BCBS~239), Bank for International Settlements, 2013. Available at \url{https://www.bis.org/publ/bcbs239.htm}.} standard sets out rigorous principles for risk data aggregation and reporting, in particular the accuracy, completeness, timeliness, and traceability of the data feeding risk reports. Together these frameworks signal an unambiguous expectation: banks must produce reliable, traceable, and auditable data on demand, and demonstrate the governance arrangements that make it trustworthy.

Despite this pressure, banks' capacity to comply varies considerably. A diagnostic survey of 44 banking professionals conducted for this project revealed three telling figures: only 30\% have a formally approved data strategy; 56.8\% routinely verify data before use, implying widespread distrust of existing quality; and 54\% experience slow or siloed access to critical information. Read together, these figures describe a sector in which the foundational conditions for regulatory compliance, a governed strategy, controlled quality, and timely access, are the exception rather than the rule. They point to a structural maturity deficit that is neither measured nor managed in any standardized way, exactly the deficit this project sets out to make visible and steerable.

\subsection{Problem Statement}

While international data maturity frameworks exist, none is simultaneously specific to the banking sector, aligned with the Tunisian regulatory environment, and operationalized through a usable assessment instrument. General purpose models describe what good data management looks like, but they do not encode the specific obligations of the BCT Circulaire N\textdegree{}2025-08 or the principles of BCBS 239. , nor do they produce a defensible, evidence backed score that a bank could present to its board or to a supervisor. As a result, Tunisian banks have no standardized means of answering a deceptively simple question: how mature is our data capability, and where must we improve to comply and to compete?

This overall problem was decomposed into four operational problems that any credible response would have to resolve. First, the absence of a structured reference framework with which to objectively evaluate the data maturity of Tunisian banks. Second, the lack of visibility on the levers of digital transformation and the business gains associated with them, which leaves investment decisions without a factual basis. Third, the difficulty institutions face in positioning themselves against data maturity standards and regulatory requirements, since no common yardstick exists. Fourth, the absence of an interactive decision-support tool that measures current maturity, identifies areas for improvement, and prioritizes the actions needed to close the gaps (the quantitative what-if simulation of business gains is scoped to future work). These four problems delimit what the project must deliver, and are addressed by the two-part solution below.

These gaps map onto the five dimensions around which the framework, introduced later, organizes the assessment: Governance (D1), Data Quality (D2), Architecture and Access (D3), Analytics and Tools (D4), and Skills and Culture (D5). Governance is rarely formalized, with most institutions lacking a board approved data strategy, assigned data ownership, and an active governance committee, leaving accountability diffuse. Data Quality is managed reactively, through manual verification rather than automated controls, so problems are detected late, by users, instead of being prevented at the source. Architecture and Access remains only partially centralized, producing slow, fragmented access and persistent data silos that undermine any single reliable version of the truth. Analytics and Tools adoption is limited, with many strategic decisions still driven by intuition and regulatory reporting that remains manual and error prone. Skills and Culture is hard to observe directly and is inferred through proxy signals such as recent data hires, training initiatives, and the involvement of business units in data decisions. The core problem this project addresses is thus the lack of a rigorous, sector specific, and regulation aligned method to measure, interpret, and steer the data maturity of Tunisian banking institutions, turning these scattered symptoms into a structured, comparable diagnosis.

\subsection{Purpose of the Study}

The purpose of this study is to design, build, and validate a means of assessing and improving data maturity in the Tunisian banking sector. The study is both analytical, diagnosing the current state of the sector and analyzing existing frameworks, and constructive, producing a solution of two connected parts. The first is a structured maturity reference framework, a hybrid model adapted to the banking context, which supplies the missing yardstick against which an institution can be objectively evaluated. The second is an interactive steering and reporting tool that operationalizes that framework for bank management: it measures the current maturity position against a framework grounded in international standards and the regulatory requirements, identifies the gaps that separate the institution from a target level derived from those standards, and does not stop at the diagnosis: it prioritizes the structuring work packages and lays them out as an improvement roadmap of concrete actions that would close the gaps and reach the target, thereby supporting decisions on the data transformation. In its final iteration the tool is delivered as a multi-tenant platform through which EY administers the instrument and each bank conducts its own assessments. The two parts are designed together, so that every concept defined in the framework is directly measurable and actionable in the tool. Their intended value is twofold. For a bank, they convert an ill defined concern about data readiness into a measured position, a set of prioritized gaps, and a defensible compliance status against the regulatory indicators. For EY, and for any advisor working with the sector, they provide a repeatable instrument that yields comparable results across institutions and successive assessments, so that progress can be tracked rather than merely asserted.

\subsection{Research Objectives}

The purpose of the study is pursued through five research objectives, which together trace the path from diagnosis to a validated, usable artifact:

\begin{enumerate}[label=(\arabic*)]
    \item Diagnose the current state of data maturity and the regulatory compliance gaps within the sector, using a diagnostic survey of banking professionals to establish an empirical baseline across the five dimensions.
    \item Analyze existing international data maturity frameworks to identify reusable components and limitations, and to justify which elements can be adapted to the banking and regulatory context of Tunisia.
    \item Design a hybrid maturity framework adapted to the banking context and to the obligations of the BCT Circulaire N\textdegree{}2025-08 and the principles of BCBS 239. The framework is structured across five weighted dimensions, twelve sub-dimensions, and 47 indicators, of which thirteen are non-skippable regulatory indicators.
    \item Construct an evidence based scoring methodology that produces an auditable and comparable maturity score on a one to five scale. The methodology incorporates an evidence cap (a score of 3 or above with no cited supporting evidence is automatically reduced to 2). It also applies a bounded skip logic (an assessor may mark a small number of non-regulatory indicators as not-applicable so they are excluded from the score, never a regulatory indicator). Finally, it incorporates a weighted aggregation into a Global Maturity Index (GMI) that is renormalized (its weights are rescaled to still sum to 100\%, so a missing dimension is not silently scored zero) when a dimension is left unanswered.
    \item Develop and validate an interactive tool, DataPilot, that operationalizes the framework for banking institutions and delivers automated scoring, visual dashboards, gap analysis, a regulatory compliance view, and an improvement roadmap. The tool is delivered as a multi-tenant platform through which EY administers the instrument and each bank conducts its own solo assessments, run by a single analyst, and group assessments, shared across a bank's departments.
\end{enumerate}

\subsection{Research Questions}

The study is structured around three research questions, each refined into sub-questions that map onto the objectives and onto the later chapters of the report:

\begin{itemize}[leftmargin=*]
    \item \textbf{RQ1.} What is the current state of data maturity in Tunisian banks, and what are the main gaps relative to international standards and regulatory requirements?
    \begin{itemize}[leftmargin=*]
        \item \textbf{RQ1a.} How do the five dimensions of Governance, Data Quality, Architecture and Access, Analytics and Tools, and Skills and Culture perform in the diagnostic baseline, and where are the most critical gaps concentrated?
        \item \textbf{RQ1b.} How compliant is the sector with the non-skippable regulatory indicators derived from the BCT Circulaire N\textdegree{}2025-08 and BCBS 239, and what regulatory exposure does the resulting compliance rate imply?
    \end{itemize}
    \item \textbf{RQ2.} How can a data maturity assessment framework be designed to fit the Tunisian banking sector, drawing on the frameworks benchmark and the survey findings?
    \begin{itemize}[leftmargin=*]
        \item \textbf{RQ2a.} Which existing data maturity frameworks (DAMA-DMBOK, DCAM, CMMI-DMM, TDWI, and Gartner) are most relevant to the Tunisian banking context, and what are their respective strengths and limitations?
        \item \textbf{RQ2b.} What dimensions, maturity levels, and scoring methodology should structure a framework adapted to Tunisian banks, based on the benchmark findings and the survey results?
    \end{itemize}
    \item \textbf{RQ3.} To what extent can an interactive data maturity steering tool support strategic decision-making and accelerate digital transformation?
    \begin{itemize}[leftmargin=*]
        \item \textbf{RQ3a.} How can the scoring methodology be operationalized so that the resulting maturity score is auditable, resistant to over-optimistic self-assessment, and comparable across institutions?
        \item \textbf{RQ3b.} How can the tool translate a maturity diagnosis into actionable guidance, through a prioritized gap analysis and a phased improvement roadmap that management can act on?
    \end{itemize}
\end{itemize}

\subsection{Approach and Boundaries of the Study}

This section delimits the study through its scope and five deliverables, its limitations, its delimitations, and its assumptions.

\subsubsection{Scope}

The study covers the design of a data maturity framework and the development of an interactive assessment tool for the Tunisian banking sector. It spans the full path from sector diagnosis and benchmark analysis through framework design, scoring methodology, and tool implementation, anchored to the obligations of the BCT Circulaire N\textdegree{}2025-08 and the principles of BCBS 239. The scope comprises five deliverables:

\begin{itemize}[leftmargin=*]
    \item a diagnostic survey of 44 banking professionals, establishing the empirical baseline of the sector across the five dimensions of the framework;
    \item a benchmark of fifteen international data maturity frameworks, from which five (DAMA-DMBOK, DCAM, CMMI-DMM, TDWI, and Gartner) were retained as sources for the hybrid design;
    \item a hybrid maturity framework structured across five weighted dimensions, twelve sub-dimensions, and 47 indicators, of which thirteen are non-skippable regulatory indicators;
    \item an evidence based scoring methodology combining an evidence cap, a bounded skip logic, and a weighted, renormalized aggregation into a Global Maturity Index on a one to five scale;
    \item DataPilot, an interactive tool that delivers the guided questionnaire, automated scoring with visual dashboards, a prioritized gap analysis, an improvement roadmap, a regulatory compliance view, and a printable export of the assessment. The tool is delivered as a multi-tenant platform with four roles (EY Admin (owner), Super Admin, Admin, Analyst), per-bank questionnaire copies from the EY master template, and department-based group assessments.
\end{itemize}

\subsubsection{Limitations}

Several limitations bound what this work establishes and how its results should be read. First, the diagnostic survey rests on 44 respondents reached during the project, so it provides indicative rather than statistically exhaustive coverage; it is not a strict probability sample, and its figures describe tendencies rather than precise population parameters. Second, the responses are self reported, carrying the risk that an organization is presented more favorably than its practices warrant; the framework mitigates this through the evidence cap built into the scoring engine, but the survey figures themselves remain subject to it. Third, the fifth dimension, skills and culture, is assessed through proxy indicators only, such as recent data hires and training initiatives, because no direct self assessment instrument for organizational culture was available; it is therefore measured less directly than the other four. Fourth, the tool was evaluated through functional testing on simulated banking data rather than a field pilot: the platform is deployed on a production instance and its behavior was verified end to end, but no real bank had produced a finalized assessment at the time of writing, so the tool is shown to be faithful to the methodology rather than validated against real institutional outcomes. Fifth, the mapping of the thirteen regulatory indicators to the BCT circular and the BCBS 239 principles is the author's own reading of those texts and was not signed off by a compliance or legal authority; expert validation of that mapping is identified as future work. Finally, the framework is calibrated to the regulatory texts in force during the project, so a subsequent regulatory change would require the regulatory indicators to be revisited.

\subsubsection{Delimitations}

The following boundaries were set deliberately. The work addresses the banking sector only, not insurance, microfinance, or other financial services, though the framework could in principle be adapted to them. It concerns the assessment and steering of data maturity, measuring the current position, exposing the gaps, and prioritizing the actions that would close them; it does not undertake the technical remediation projects an assessment may recommend, which remain the institution's responsibility. The scoring model is auditable and comparable rather than predictive: it does not forecast the financial return of a data transformation, and the quantitative what-if simulation of business gains is left to future work. The project is a design and construction effort rather than an empirical validation study, so it includes no controlled comparison of outcomes across institutions.

\subsubsection{Assumptions}

The work rests on a few explicit assumptions, stated so the reader can weigh the results accordingly. The survey respondents are assumed to have answered in good faith, reflecting their institutions' actual practices within the limits of self reporting. The regulatory texts in force during the project, principally the BCT Circulaire N\textdegree{}2025-08 and the BCBS 239 principles, are assumed to remain the relevant compliance reference for its duration, and their obligations to be correctly interpreted in the regulatory indicators. The five international frameworks retained from the benchmark are assumed to represent the established state of practice and thus a sound basis for a hybrid design. Finally, the seeded demonstration data used to exercise the tool is assumed representative enough of a realistic assessment to test the fidelity of the scoring, even though it does not stand in for a real institutional measurement.

\subsection{Contributions}

This report presents three artifacts: the benchmark and hybrid framework, the scoring methodology with its rubrics, and the DataPilot tool, all designed, written, and implemented by the author during the six month internship under the guidance of the academic and institution supervisors. They constitute three contributions. The first is conceptual: a hybrid data maturity framework tailored to the banking sector, structured across five dimensions and 47 evidence based indicators distributed over twelve sub-dimensions. Thirteen of these are non-skippable regulatory indicators explicitly aligned with the obligations of the BCT Circulaire N\textdegree{}2025-08 and the principles of BCBS 239. The second is methodological: a rigorous evidence based scoring methodology with defined anti-gaming rules, a skip logic mechanism, and a structured weighting system that produces an auditable global maturity score on a one to five scale. The third is practical: DataPilot, an interactive assessment tool that operationalizes the framework, providing bank management with automated scoring, visual dashboards, gap analysis, a regulatory compliance view, and an improvement roadmap across the six screens of the assessment workflow. It is delivered as a multi-tenant platform with four roles (EY Admin (owner), Super Admin, Admin, Analyst), invite-only onboarding, bank-scoped tenant isolation, department-based group assessments, and per-bank questionnaire copies of the EY master template.

\subsection{Structure}

The remainder of this report is organized as follows. Chapter 2 presents the background and related studies, covering the banking and regulatory environment, the concept of data maturity, and an analysis of the five international frameworks, and ends with the research gap that motivates the design. Chapter 3 describes the methodology, combining Design Science Research with an iterative, Agile inspired development approach, together with the project phases and planning, the stakeholder analysis, the research materials, and the tools. Chapter 4, the core of the report, presents the four iterations through which the artifacts, the framework, the scoring engine, the DataPilot tool, and the multi-tenant platform, were designed, validated, implemented, and evaluated, detailing the dimensions, the scoring rules, the software architecture, and the use case specifications. Chapter 5 presents the results, including the diagnostic maturity profile and a worked example of the Global Maturity Index, and answers the research questions. Chapter 6 discusses the findings, their theoretical and practical implications, and the threats to validity with their mitigations. Chapter 7 concludes and outlines future work.

\subsection{Conclusion}

This chapter introduced the host company, the context and problem motivating the study, its purpose, objectives, and research questions, the boundaries of the work, and its contributions. The next chapter reviews the background and related studies that ground the design of the framework.
